% !TeX program = pdflatex
% !BIB program = biber
\documentclass[11pt,a4paper]{article}
\usepackage[T1]{fontenc}
\usepackage[utf8]{inputenc}
\usepackage[french]{babel}
\usepackage{lmodern,microtype}
\usepackage[margin=2.1cm]{geometry}
\usepackage{amsmath}
\usepackage{booktabs,tabularx,array}
\usepackage{enumitem}
\usepackage{hyperref}
\usepackage{xcolor}
\usepackage{float}
\usepackage{svg}
\usepackage[backend=biber,style=numeric,sorting=none,maxnames=2,giveninits=true,doi=true,url=false,isbn=false]{biblatex}
\addbibresource{ref.bib}
\hypersetup{colorlinks=true,linkcolor=blue!50!black,citecolor=blue!50!black,urlcolor=blue!50!black}
\setlist[itemize]{nosep,leftmargin=*}
\setlist[enumerate]{nosep,leftmargin=*}
\renewcommand{\arraystretch}{1.17}
\newcommand{\todo}[1]{\textcolor{blue!45!black}{\emph{#1}}}
\title{\textbf{Perception visuelle par caméras fisheye embarquées}}
\author{Sala Raphaël, CHOUMAN Ahmad }
\date{DATE}
\begin{document}
\maketitle
% \noindent\textbf{Nature du document.} Proposition de plan et de références pour un rapport final d'environ \textbf{15 pages bibliographie comprise}. Il ne s'agit pas de résultats déjà démontrés. Les choix techniques et protocoles sont à vérifier lors des essais. On privilégie la compréhension visuelle des méthodes, l'informatique expérimentale et les preuves de reproductibilité.


\begin{abstract}
Les caméras fisheye offrent un champ de vision étendu. Leur géométrie de projection, différente de celle des caméras perspectives, complique toutefois l'utilisation directe des méthodes classiques de vision. Ce projet étudie la calibration de ces capteurs et l'influence de différentes représentations de l'image sur deux tâches complémentaires : la détection et le suivi de personnes dans un flux vidéo, et l'estimation d'une carte de profondeur dense à partir d'un banc stéréo fisheye. À terme, leur association vise à localiser les personnes suivies dans le repère du robot. Ce rapport présente les fondements géométriques, les approches retenues dans la littérature et les critères envisagés pour leur évaluation.
\end{abstract}


\textcolor{red}{le sommaire est juste là pour vérifier la pagination et la structure, il sera retiré dans le rapport final.}

\tableofcontents


% \section*{Objectif et périmètre}
% \noindent \textbf{Question du projet :} comment exploiter des images de caméras fisheye avec des modèles de vision préentraînés sur images perspectives pour (i) détecter et suivre des personnes, (ii) calculer une profondeur dense à partir d'un banc fisheye stéréo, puis (iii) réunir les deux sorties pour connaître la position 3D des personnes dans le repère du robot ?
% \begin{itemize}
% \item Entrées : flux fisheye calibrés, et paires stéréo synchronisées quand disponibles.
% \item Branche 1 : reprojection cylindrique ou ERP $\to$ détecteur préentraîné $\to$ tracker multi-objets $\to$ identités/trajectoires 2D.
% \item Branche 2 : calibration relative $\to$ rectification stéréo adaptée au banc $\to$ correspondances denses $\to$ profondeur métrique, avec masque valide.
% \item Fusion envisagée : lier à chaque instance suivie un point/une zone 3D et construire une trajectoire en coordonnées robot.
% \end{itemize}
% \noindent\textbf{Deux précisions indispensables.} Reprojeter sur un cylindre ou une sphère \emph{ne transforme pas} toute une image fisheye en une seule image perspective classique. Et une carte de profondeur n'est disponible que là où les deux caméras voient la même scène avec une géométrie exploitable : \emph{champ total et champ stéréo commun ne coïncident pas nécessairement}.


\section{Introduction}

Les caméras fisheye occupent une place importante en vision par ordinateur dès qu'un large champ de vue est recherché. On les rencontre en robotique, en surveillance, en imagerie médicale, en réalité virtuelle et dans de nombreux autres domaines \textcolor{red}{REF}.
Leur intérêt principal est de capturer un champ de vision beaucoup plus large qu'une caméra perspective classique. Toutefois les objets observés loin de l'axe optique apparaissent fortement déformés et les méthodes de vision classiques, souvent conçues pour des images rectilignes, ne s'appliquent pas directement.

\textcolor{red}{Je réfléchis toujours à insérer un schema explicatif + voir si on aura besoin des ces notations. Je laisse tel quel pour l'instant.}
Pour décrire cette différence, on considère un point 3D exprimé dans le repère de la caméra de coordonnées $X_c=(X,Y,Z)^\top$. La direction de ce point peut être décrite par deux angles : l'angle polaire $\theta$ qui mesure son écart à l'axe optique et l'azimut $\varphi$ qui indique sa direction autour de cet axe. Dans un modèle fisheye central, l'idée essentielle est que l'on ne projette pas directement le point 3D sur un plan image comme dans le modèle perspective classique. On commence d'abord par relier l'angle $\theta$ à une distance radiale dans l'image, notée $q(\theta)$.
Cette fonction $q$ décrit donc la manière dont une direction de l'espace est reportée dans l'image. Une fois cette distance radiale obtenue, on la répartit selon la direction angulaire $\varphi$, puis on convertit la position obtenue en coordonnées pixels grâce aux paramètres intrinsèques de la caméra, en particulier les focales $f_x,f_y$ et le point principal $(c_x,c_y)$. Ainsi, le rôle de $q$ est central, c'est elle qui détermine comment les écarts angulaires dans la scène se traduisent en distances dans l'image et donc la nature même de la déformation fisheye.


Plusieurs modèles de projection se distinguent par la loi radiale $q(\theta)$ qui associe à chaque direction de visée une position dans l'image (figure~\ref{fig:comparaison-projections}). Le modèle rectiligne utilise $q(\theta)=\tan\theta$ et permet que les droites de l'espace restent des droites dans l'image. En revanche, cette fonction diverge lorsque $\theta$ approche $90^\circ$, ce qui empêche de représenter sur un seul plan perspective un champ atteignant $180^\circ$. Les projections fisheye classiques répartissent autrement les directions observées. Le modèle équidistant fait correspondre des écarts angulaires égaux à des écarts radiaux égaux, tandis que le modèle équisolide conserve les rapports d'angles solides sous forme de rapports d'aires dans l'image. Le modèle stéréographique préserve localement les angles, alors que la projection orthographique correspond à la projection d'un hémisphère sur un plan~\cite{schneider2009,yogamani2019}. Ces propriétés ont des conséquences sur la déformation des objets et la distribution des détails dans l'image. Les domaines de validité diffèrent également, par exemple la loi orthographique $q(\theta)=\sin\theta$ cesse d'être injective au-delà de $90^\circ$.

Ces modèles sont toutefois des descriptions idéales. La projection d'une optique réelle s'en écarte en raison des imperfections de fabrication et d'alignement. La calibration consiste alors à estimer un modèle capable de reproduire les positions effectivement mesurées dans l'image. Kannala et Brandt proposent notamment une loi radiale polynomiale qui généralise la projection équidistante et dont les coefficients sont ajustés à partir d'observations de calibration~\cite{kannala2006}. D'autres familles reposent sur une construction géométrique différente. Le modèle unifié de Mei et Rives fait intervenir une projection sur une sphère~\cite{mei2007}, tandis que le modèle Double Sphere utilise deux sphères et permet un calcul efficace de la projection inverse~\cite{usenko2018}. Aucun modèle n'est optimal pour toutes les optiques. Leur pertinence se juge notamment par la précision de reprojection, le domaine angulaire couvert et le coût des transformations. Une calibration fiable est ainsi nécessaire pour relier les pixels aux directions de l'espace et pour réaliser ensuite des reprojections ou des traitements stéréoscopiques cohérents~\cite{schneider2009,yogamani2019}. \textcolor{red}{J'essaierai de m'y intéresser plus pour expliquer pourquoi on choisit l'un oplutôt que l'autre dans la section Projet. Ca dépendra de commnent la calibration banc stéréo est présenté afin de pas faire redondant ou de trop technique.}

\begin{figure}[H]
  \centering
  \includesvg[width=\linewidth]{img/fisheye_modeles_projection_1x5.svg}
  \caption{Comparaison de cinq lois de projection appliquées à un même damier en 3D. La fonction $q(\theta)$ est indiquée pour chaque modèle. \textcolor{red}{TODO : grossir texte image}}
  \label{fig:comparaison-projections}
\end{figure}


La suite de ce rapport s'organise autour de quatre parties. Nous étudions d'abord les principaux prétraitements des images fisheye, en particulier les reprojections cylindrique et équirectangulaire, ainsi que leurs limites. Nous présentons ensuite les approches de détection et de suivi de personnes, avant d'aborder la calibration d'un banc stéréo fisheye et les méthodes d'estimation de profondeur dense. Enfin, nous revenons sur les choix méthodologiques et les avancées du projet, avec pour objectif à terme de combiner le suivi et la profondeur afin de localiser les personnes dans l'espace.




\section{Prétraitements (3 pages)}

\textcolor{red}{Principe du pré-traitement, justification et bref énumération des diffrentes méthodes avant de se concentrer sur cylindrique et ERP, analyse qualitative (car quantitative c'est pas vraiment possible je crois)}

\textcolor{red}{Ahmad (début de redac) :}

Les modèles de vision classiques (comme les réseaux convolutionnels ou les transformateurs) ont été conçus et entraînés pour des images standards prises par des caméras classiques. Ils s'attendent à ce qu'une personne ou un objet garde la même forme générale, peu importe où il se trouve dans l'image. 

Avec un objectif fisheye, ce n'est plus du tout le cas. Plus un objet s'éloigne du centre pour aller vers les bords, plus il est écrasé, courbé et déformé. Les modèles standards perdent donc vite leurs repères si on leur donne une image fisheye brute. Pour contourner ce problème sans devoir réentraîner un modèle de zéro sur des millions d'images fisheye annotées, la solution la plus simple consiste à « déplier » ou redresser géométriquement l'image sur une surface virtuelle intermédiaire. 

Nous nous concentrons ici sur les deux représentations les plus courantes : la projection cylindrique et la projection équirectangulaire (ERP).

\subsection{Cadre géométrique et principe du redressement inverse}

Pour passer d'une image fisheye à une image dépliée, on commence par repérer chaque direction de l'espace. En prenant le centre optique de la caméra comme origine $\mathcal{R}_c = (O_c, \mathbf{X}_c, \mathbf{Y}_c, \mathbf{Z}_c)$ avec l'axe $\mathbf{Z}_c$ pointé tout droit, la direction de n'importe quel point 3D $\mathbf{X} = (X, Y, Z)^\top$ est définie par deux angles :
\begin{itemize}
    \item L'angle $\theta \in [0, \pi]$, qui mesure à quel point la direction s'écarte de l'axe central (l'axe optique).
    \item L'angle $\varphi \in (-\pi, \pi]$, qui indique l'orientation tout autour de cet axe.
\end{itemize}
Mathématiquement :
\begin{equation}
\theta = \arccos\left(\frac{Z}{\|\mathbf{X}\|_2}\right), \quad \varphi = \operatorname{atan2}(Y, X).
\end{equation}

Une caméra fisheye réelle suit généralement un modèle géométrique calibré, comme celui de Kannala-Brandt[cite: 1]. Dans ce modèle, l'angle $\theta$ détermine à quelle distance $r(\theta)$ du centre de l'image le point va atterrir :
\begin{equation}
r(\theta) = k_1\theta + k_2\theta^3 + k_3\theta^5 + k_4\theta^7.
\end{equation}
Les coefficients $k_i$ et les paramètres internes (focales $f_x, f_y$ et centre $(c_x, c_y)$) permettent ensuite de trouver les coordonnées pixels réelles $(u_f, v_f)$ sur le capteur :
\begin{equation}
\begin{pmatrix} u_f \\ v_f \end{pmatrix} = \begin{pmatrix} f_x & 0 \\ 0 & f_y \end{pmatrix} \begin{pmatrix} r(\theta)\cos\varphi \\ r(\theta)\sin\varphi \end{pmatrix} + \begin{pmatrix} c_x \\ c_y \end{pmatrix}.
\end{equation}

\paragraph{Pourquoi utilise-t-on le redressement inverse ?}
Si on projetait directement chaque pixel du fisheye vers l'image dépliée, on se retrouverait avec des « trous » et des pixels noirs dans l'image finale. On procède donc en marche arrière (\emph{backward mapping}) : pour chaque pixel de l'image de destination, on calcule la direction 3D vers laquelle il regarde, on cherche où cette direction atterrit dans l'image fisheye de base, puis on interpole la couleur avec les pixels voisins.

\subsection{Reprojection cylindrique}

L'idée est simple : on entoure la caméra d'un cylindre virtuel vertical, puis on déroule ce cylindre à plat comme une feuille de papier[cite: 1]. En robotique mobile terrestre, c'est particulièrement pratique car le robot et les piétons évoluent sur le même sol horizontal[cite: 1].

\paragraph{Formulation mathématique.}
Sur la feuille déroulée, un point est repéré par ses coordonnées planes $(x_{\text{cyl}}, y_{\text{cyl}})$. Si le cylindre a un rayon $R$, ces coordonnées sont reliées à l'angle horizontal $\psi$ (l'azimut) et à la hauteur $Y$ par :
\begin{equation}
x_{\text{cyl}} = R\,\psi, \quad y_{\text{cyl}} = R\,\frac{Y}{\sqrt{X^2 + Z^2}} = R\,\tan \phi_y,
\end{equation}
où $\psi = \operatorname{atan2}(X, Z)$ désigne l'angle horizontal par rapport à l'axe de visée, et $\phi_y$ l'angle vers le haut ou vers le bas.

\textcolor{red}{[TODO : finir d'expliquer l'avantage pour les piétons debouts et l'étirement vers le haut + ajouter images]}

\subsection{Projection équirectangulaire (ERP)}

\paragraph{Formulation mathématique.}
\textcolor{red}{[TODO : expliquer l'ERP puis pourquoi ça étire les pôles. À finir plus tard.]}

\textcolor{red}{images}


\section{Détection et suivi (4 pages)}
\subsection*{Détection}

\textcolor{red}{Décrire brièvement la logique d'un détecteur préentraîné (YOLOv8 / YOLO11 ; RT-DETR~\cite{zhao2024} en comparaison si justifiée).
Cout computationnel, vitesse, précision, classes disponibles, AP50, mAP50--95. Montrer les limites de la détection sur des images fisheye brutes et les améliorations obtenues par reprojection cylindrique ou ERP. Mentionner brièvement les alternatives (FishEye8K~\cite{plauttutorial}, WoodScape~\cite{bourkefish2pano}) et leurs limites (données synthétiques, classes limitées, performances moindres). Montrer que la détection est un prérequis pour le suivi et la localisation 3D.}

\subsection*{Tracking}
\textcolor{red}{Principe du tracking, filtre de Kalmann, occlusion,cout computationnel}

Expliquer (brievement) les différentes approches/modèles la chronologie, les nouveautés, les performances \textbf{SORT}~\cite{bewley2016} (Kalman + IoU + affectation), \textbf{Deep SORT}~\cite{wojke2017} (apparence), \textbf{ByteTrack}~\cite{zhang2022bytetrack} (détections de confiance basse), \textbf{BoT-SORT}~\cite{aharon2022} (mouvement/apparence). Puis deux publications spécialisées au \emph{suivi fisheye} particulièrement utiles : \textbf{FORT}~\cite{odic2023fort} et suivi de boîtes orientées \textbf{Lam (BMVC 2024)}~\cite{lam2024}. 


\section{Banc stéréo fisheye et estimation dense (4 pages)}
\subsection*{Calibration relative et rectification}
\noindent Présenter la chaîne : deux vues $\to$ poses calibrées $\to$ rectification épipolaire adaptée à la géométrie $\to$ correspondances. Une projection panoramique \emph{seule} n'est pas équivalente à une rectification stéréo. Références à présenter dans un ordre historique lisible : Abraham--Förstner (2005)~\cite{abraham2005} ; Ohashi (2016/2017)~\cite{ohashi2016,ohashi2017} ; Schneider (2016)~\cite{schneider2016}. Comme mise en perspective plus moderne, mentionner les stratégies de rectification d'\textbf{OmniVidar}~\cite{xie2023} et de projection de Cassini dans \textbf{OmniStereo}~\cite{deng2025}. Les méthodes multi-fisheye ne sont pas toutes transférables à deux caméras ; préciser le nombre et la disposition des capteurs des articles.
\noindent \textbf{Figure requise :} deux centres de caméras, rayons vers un point 3D et plan/ligne épipolaire ; puis une paire réellement rectifiée avec lignes guides.



\subsection*{Cartes denses et évaluation}
\textcolor{red}{decrire ce que sont les cartes de profondeur, comment elles sont faites : creuses VS denses. On se concentre sur la derniere car [??]}

Classique : \textbf{StereoSGBM}~\cite{hirschmuller2008}, puis modèles préentraînés de correspondance stéréo \textbf{RAFT-Stereo}~\cite{lipson2021}, \textbf{FoundationStereo}~\cite{wen2025} (et éventuellement \textbf{Fast-FoundationStereo}, 2026~\cite{wen2026fast}). En état de l'art stéréo panoramique, citer \textbf{OmniMVS}~\cite{won2019}, \textbf{OmniVidar}~\cite{xie2023}, \textbf{OmniStereo}~\cite{deng2025}, \textbf{DFI-OmniStereo}~\cite{endres2025} et \textbf{MDP-Omni}~\cite{son2025} en les situant selon leur rig et leur projection. \textbf{WideDepth}~\cite{indyk2026} est un benchmark pertinent pour une première validation ; ses images de test fisheye sont produites à partir de scans laser de scènes réelles : ne pas les confondre avec une séquence brute capturée par vos caméras. HELVIPAD~\cite{zayene2025} est utile pour les images ERP, mais provient d'un autre type de configuration (deux caméras omnidirectionnelles superposées).
% sugestion assez détaillée de chatGPT
% \noindent \textbf{Figure requise :} images gauche/droite, disparité, profondeur colorisée, masque valide et (si disponible) carte d'erreur vs référence. Distinguer profondeur axiale, distance au centre et disparité ; ne pas appliquer aveuglément $Z=fB/d$ à des coordonnées sphériques non rectifiées.


\section{Projet}

Parler de notre projet en particulier, des différentes avancées. De la pipeline pour aborder la détection. Je pense le but final c'est d'arriver à traquer dans un espace 3d, un groupe de personnes qui marche et de reussir à savori où elles se trouvent. De manière détaillée ça serait : 

Pour chaque personne détectée et suivie, remapper sa zone en coordonnées de la carte de profondeur, extraire une estimation robuste de distance (avec masque valide), reconstruire un point 3D et le transformer dans le repère robot. Le point 3D estimé sera un \emph{point représentatif} et non une boîte 3D complète, sauf extension décidée explicitement. Une médiane des valeurs de profondeur dans une région centrale est une première approche pratique, mais peut être biaisée par le fond ; mentionner le besoin éventuel de segmentation ou d'un point d'appui pertinent. Les ID sont définis par le tracker, pas par la carte de profondeur.
\noindent \textbf{Figure requise :} image avec bbox personne et ID, carte de profondeur correspondante, même ID associé à un point $(X,Y,Z)$ dans le repère du robot.

\begingroup
\small
\printbibliography[title={Bibliographie citée dans ce plan}]
\endgroup
\end{document}
